% 页面与段落
\onehalfspacing
\setlength{\parindent}{2em}
\setlength{\parskip}{0pt}
\setlength{\headheight}{15pt}
\setlist{nosep,leftmargin=2em}

% 标题
\ctexset{
  section={format=\centering\heiti\zihao{3},beforeskip=1.4ex,afterskip=1.0ex},
  subsection={format=\heiti\zihao{4},beforeskip=1.2ex,afterskip=0.7ex},
  subsubsection={format=\heiti\zihao{-4},beforeskip=1.0ex,afterskip=0.5ex}
}
\setcounter{secnumdepth}{3}
\setcounter{tocdepth}{2}
\numberwithin{equation}{section}

% 页眉页脚
\pagestyle{fancy}
\fancyhf{}
\fancyhead[C]{2026 年高教社杯全国大学生数学建模竞赛 B 题}
\fancyfoot[C]{\thepage}
\renewcommand{\headrulewidth}{0.4pt}
\renewcommand{\footrulewidth}{0pt}

% 图表与链接
\captionsetup{font=small,labelsep=quad}
\hypersetup{
  colorlinks=true,
  linkcolor=black,
  citecolor=black,
  urlcolor=NavyBlue,
  pdfauthor={参赛队伍},
  pdftitle={无线电干扰源的快速自动定位与清除}
}

% 代码清单
\lstdefinestyle{modeling}{
  language=Python,
  basicstyle=\ttfamily\small,
  keywordstyle=\color{NavyBlue}\bfseries,
  commentstyle=\color{OliveGreen},
  stringstyle=\color{BrickRed},
  numbers=left,
  numberstyle=\tiny\color{gray},
  frame=single,
  breaklines=true,
  showstringspaces=false,
  columns=fullflexible
}
\lstset{style=modeling}
